\documentclass[11pt]{article}
\usepackage[margin=2.5cm]{geometry}
\usepackage{amsmath,amssymb,amsthm}
\usepackage[hidelinks]{hyperref}
\newtheorem{theorem}{Theorem}[section]
\newtheorem{lemma}[theorem]{Lemma}
\newtheorem{proposition}[theorem]{Proposition}
\newtheorem{corollary}[theorem]{Corollary}
\theoremstyle{definition}
\newtheorem{definition}[theorem]{Definition}
\newtheorem{remark}[theorem]{Remark}
\newcommand{\Z}{\mathbb{Z}}
\newcommand{\PP}{\mathbb{P}}
\newcommand{\B}{\mathcal{B}}
\newcommand{\dB}{d_{\mathcal{B}}}
\newcommand{\DB}{D_{\mathcal{B}}}
\newcommand{\R}{R}
\renewcommand{\le}{\leqslant}
\renewcommand{\ge}{\geqslant}
\renewcommand{\leq}{\leqslant}
\renewcommand{\geq}{\geqslant}
\author{Problem 3 team}

\title{Problem 3 (Bulgarian solitaire), Cell 3:\\ a general upper bound}
\date{}
\begin{document}
\maketitle

\paragraph{Official setting.}
\begin{quote}
A partition of a positive integer $n$ is a weakly decreasing sequence $\lambda=(\lambda_1,\dots,\lambda_s)$ of
positive integers with $\lambda_1+\cdots+\lambda_s=n$.  The shift $\B(\lambda)$ is the partition whose parts are the
positive numbers among $\lambda_1-1,\dots,\lambda_s-1$ together with one extra part equal to $s$.  Call $\lambda$
cyclic if $\B^i(\lambda)=\lambda$ for some $i\ge1$, and let $\dB(\lambda)=\min\{\,i\ge0:\B^i(\lambda)\text{ is cyclic}\,\}$,
$\DB(n)=\max\{\,\dB(\lambda):\lambda\text{ is a partition of }n\,\}$.  Write $T_k=k(k+1)/2$ and
$\delta_k=(k,k-1,\dots,2,1)$.  Every $n\ge1$ satisfies $T_{k-1}<n\le T_k$ for exactly one $k$, which we call the
rank of $n$.
\end{quote}
Section~\ref{sec:defs} restates these definitions in an equivalent form, proves that $\dB$ is well defined and
that the rank exists and is unique (Lemma~\ref{lem:rank}, where we also write $n=T_{k-1}+r$, $1\le r\le k$), and
fixes further notation.

\paragraph{Official statement of the cell.}
\begin{quote}
Prove that for every $k\ge4$ and every non-triangular $n$ with $T_{k-1}<n<T_k$, $\DB(n)\le k^2-2k-1$, and
determine $\DB(T_k-1)$ exactly.  Determine also, for that $n$, which partitions attain the maximum.
\end{quote}

\paragraph{Status.} \textbf{Partial:} value and both bounds fully proved for the stated range; the complete set of extremal partitions is given only by an exact dynamical criterion (closed form open).
\begin{itemize}
\item The bound and the value are proved in full.  The case $k=4$ uses an exhaustive computation over the $67$
partitions of $7,8,9$, with the code in Appendix~\ref{app:code}.
\item The partitions attaining the maximum are determined \emph{exactly} by the criterion of
Theorem~\ref{thm:c3ext}.  We also give an explicit extremal family and a proved direct necessary condition.
\item Qualification: this is a dynamical criterion, not a closed-form description of all extremal partitions;
that description is open (Section~\ref{sec:c3direct}).  If a closed form is required, this sub-part is partial.
\end{itemize}

\paragraph{Answer and where it is proved.}
\begin{itemize}
\item For $k\ge4$ and $T_{k-1}<n<T_k$: $\DB(n)\le k^2-2k-1$ (Theorem~\ref{thm:c3upper}, \cite[Thm.~4.4(1)]{GH}).
\item $\DB(T_k-1)=k^2-2k-1$ for every $k\ge4$ (Theorem~\ref{thm:c3lower}), attained by the explicit family
$\lambda^{\langle k\rangle}=(k-1,\,k-2,\,k-2,\,k-3,\dots,2,\,1,\,1)$ ($\lambda_1=k-1$, $\lambda_2=k-2$,
$\lambda_i=k+1-i$ for $3\le i\le k$, $\lambda_{k+1}=1$).
\item Which partitions attain it.  For $k=4$: only $(3,2,2,1,1)$.  For $k\ge5$: exactly those $\lambda\vdash T_k-1$ for which
$\B^{k^2-2k-3}(\lambda)$ has at least $k+1$ parts; equivalently, the part counts along the orbit contain the ramps
$(k-1,k,\dots,k,k+1)$ at the positions $jk,\dots,jk+j+1$, $j=1,\dots,k-3$ (Theorem~\ref{thm:c3ext}).  Every extremal
partition ($k\ge5$) satisfies the direct conditions $\lambda_1\le k+1$ and $\#\{m:\lambda_m\ge i\}\ge k+3-2i$ for all
$i\ge1$ (Proposition~\ref{prop:c3direct}).
\end{itemize}
The proof of the bound reproduces, completed in detail, Griggs--Ho \cite[Lemma~4.3, Thm.~4.4]{GH}.  Their proof of
the tail lemma is only sketched there; we give a different complete proof.  Sections~\ref{sec:defs}--\ref{sec:tail}
are shared with Cell 2 (and, up to Section~\ref{sec:ramps}, with Cells 4 and 5), and are repeated so that this
document is self-contained.

\section{Setting and notation}\label{sec:defs}

Throughout, $T_m=m(m+1)/2$ for integers $m\ge 0$, and $[P]\in\{0,1\}$ denotes the truth value
of a statement $P$.
A \emph{partition} $\lambda$ of $n\ge 1$ (written $\lambda\vdash n$) is a finite non-increasing
sequence of positive integers $\lambda_1\ge\lambda_2\ge\dots\ge\lambda_\ell$ with sum $n$;
$\ell=\ell(\lambda)$ is its number of parts, and we put $\lambda_i=0$ for $i>\ell(\lambda)$.
Partitions are identified with multisets of positive integers.

\begin{definition}[Bulgarian solitaire]
For $\lambda\vdash n$ let $\B(\lambda)$ be the partition whose parts are $\ell(\lambda)$ together
with the numbers $\lambda_i-1$ for those $i\le\ell(\lambda)$ with $\lambda_i\ge 2$.
Since $\ell(\lambda)+\sum_{i}(\lambda_i-1)=n$, we have $\B(\lambda)\vdash n$.
Put $\B^0=\mathrm{id}$ and $\B^{i+1}=\B\circ\B^{i}$.
A partition $\mu$ is \emph{cyclic} if $\B^P(\mu)=\mu$ for some $P\ge 1$.
For $\lambda\vdash n$, $\dB(\lambda)$ is the least $i\ge 0$ such that $\B^i(\lambda)$ is cyclic,
and $\DB(n)=\max\{\dB(\lambda):\lambda\vdash n\}$.
A partition $\lambda\vdash n$ with $\dB(\lambda)=\DB(n)$ is called \emph{extremal}.
\end{definition}

$\dB(\lambda)$ is well defined: $n$ has finitely many partitions, so $\B^a(\lambda)=\B^{a+P}(\lambda)$
for some $a\ge 0$, $P\ge 1$, and then $\B^a(\lambda)$ is cyclic.
If $\mu$ is cyclic with $\B^P(\mu)=\mu$, then $\B^P(\B(\mu))=\B(\mu)$, so $\B(\mu)$ is cyclic.
Hence $\B^i(\lambda)$ is cyclic for every $i\ge\dB(\lambda)$ and for no $i<\dB(\lambda)$.

\begin{lemma}[rank]\label{lem:rank}
Every $n\ge1$ can be written in exactly one way as $n=T_{k-1}+r$ with integers $k\ge1$ and
$1\le r\le k$.  We call $k$ the \emph{rank} of $n$; $n$ is triangular iff $r=k$ (then $n=T_k$).
\end{lemma}
\begin{proof}
$T_0=0<T_1<T_2<\cdots$ is strictly increasing and unbounded, so there is exactly one $k\ge 1$ with
$T_{k-1}<n\le T_k$; then $r=n-T_{k-1}$ satisfies $1\le r\le T_k-T_{k-1}=k$.
Conversely $n=T_{k-1}+r$ with $1\le r\le k$ forces $T_{k-1}<n\le T_k$, i.e.\ this $k$.
If $n=T_m$ ($m\ge1$), then $T_{m-1}<n\le T_m$, so $k=m$ and $r=T_m-T_{m-1}=m=k$; conversely $r=k$ gives
$n=T_k$.
\end{proof}

\section{Diagrams, diagonals, rotation, and Brandt's theorem}\label{sec:brandt}

This section follows Griggs--Ho \cite[\S2]{GH} (the idea goes back to Etienne \cite{Et});
we make their ``left shifting decreases the diagonal index'' precise with an explicit
potential $\Phi$.

Let $\PP=\Z_{>0}\times\Z_{>0}$; the \emph{cell} $(i,j)\in\PP$ has \emph{row} $i$ and \emph{column} $j$.
The \emph{diagram} of $\lambda$ is
\[Y(\lambda)=\{(i,j)\in\PP:\ 1\le j\le\ell(\lambda),\ 1\le i\le\lambda_j\},\]
so column $j$ has $\lambda_j$ cells (columns are parts, as in \cite{GH}).
$Y(\lambda)$ determines $\lambda$, since $\lambda_j=\#\{i:(i,j)\in Y(\lambda)\}$.
Row~1 of $Y(\lambda)$ is $\{(1,j):j\le\ell(\lambda)\}$ and column~1 has $\lambda_1$ cells; in particular
\begin{equation}\label{eq:Y1}
(1,w+1)\in Y(\lambda)\ \Longrightarrow\ (1,w)\in Y(\lambda)\qquad(w\ge1).
\end{equation}
For a finite sequence $a=(a_1,\dots,a_m)$ of non-negative integers let
$X_a=\{(i,j):1\le j\le m,\ 1\le i\le a_j\}$.  If $a$ is non-increasing, then $X_a=Y(\nu)$ where $\nu$ is the
partition formed by the non-zero entries of $a$.

For $w\ge 1$ the \emph{$w$-th diagonal} is $D_w=\{(i,j)\in\PP: i+j=w+1\}$; it has the $w$ cells
$d_w(j)=(w+1-j,\,j)$, $1\le j\le w$ ($d_w(j)$ is the cell of $D_w$ in column $j$).
$\PP$ is the disjoint union of the $D_w$.  Put $\Delta_0=\emptyset$ and
$\Delta_m=D_1\cup\dots\cup D_m$ for $m\ge 1$.  Column $j$ of $\Delta_m$ consists of the cells $(i,j)$ with
$i\le m+1-j$, so $\Delta_m=Y(\delta_m)$ with $\delta_m=(m,m-1,\dots,1)$, and $|\Delta_m|=T_m$.

Define $\rho:\PP\to\PP$ by $\rho(i,j)=(i-1,j+1)$ if $i\ge2$ and $\rho(1,j)=(j,1)$.

\begin{lemma}\label{lem:rho}
(a) $\rho(d_w(j))=d_w(j+1)$ for $1\le j<w$ and $\rho(d_w(w))=d_w(1)$.  Hence $\rho$ maps each $D_w$
bijectively onto itself, $\rho$ is a bijection of $\PP$, and for $m\ge 0$,
$\rho^m(d_w(j))=d_w(j')$ where $1\le j'\le w$ and $j'\equiv j+m\pmod w$.  In particular $\rho^w$ is the
identity on $D_w$.
(b) $\rho$ preserves the quantity $i+j$ of a cell.
\end{lemma}
\begin{proof}
(a) For $j<w$ the cell $d_w(j)=(w+1-j,j)$ has row $w+1-j\ge2$, so $\rho(d_w(j))=(w-j,j+1)=d_w(j+1)$;
and $\rho(d_w(w))=\rho(1,w)=(w,1)=d_w(1)$.  So $\rho$ acts on $D_w$ as the cyclic shift $j\mapsto j+1$
of column indices modulo $w$; the rest follows.  (b) $(i-1)+(j+1)=i+j$ and $j+1=1+j$.
\end{proof}

\begin{lemma}[one move]\label{lem:move}
Let $\lambda\vdash n$, $s=\ell(\lambda)$ and $a=(s,\lambda_1-1,\lambda_2-1,\dots,\lambda_s-1)$.
Then $\rho(Y(\lambda))=X_a$, and $Y(\B(\lambda))=X_{a^*}$, where $a^*$ is the non-increasing rearrangement
of $a$.  If $s\ge\lambda_1-1$, then $a$ is non-increasing and $Y(\B(\lambda))=\rho(Y(\lambda))$.
\end{lemma}
\begin{proof}
The cells of row~1 of $Y(\lambda)$ are $(1,j)$, $1\le j\le s$; $\rho$ sends them to $(j,1)$, i.e.\ exactly
onto the cells $(1,1),\dots,(s,1)$ of column~1.  A cell $(i,j)\in Y(\lambda)$ with $i\ge 2$ satisfies
$2\le i\le\lambda_j$ and is sent to $(i-1,j+1)$; for fixed $j$ these images are exactly the cells
$(1,j+1),\dots,(\lambda_j-1,j+1)$.  Hence $\rho(Y(\lambda))=X_a$.  The parts of $\B(\lambda)$ are the non-zero
entries of $a$, so $Y(\B(\lambda))=X_{a^*}$.  The entries of $a$ after the first are non-increasing because
$\lambda$ is; so $a$ is non-increasing iff $s\ge\lambda_1-1$, and then $a^*=a$.
\end{proof}

For a finite set $X\subseteq\PP$ put $\Phi(X)=\sum_{(i,j)\in X}(i+j)$.

\begin{lemma}\label{lem:rearr}
For every finite sequence $a$ of non-negative integers, $\Phi(X_{a^*})\le\Phi(X_a)$, with equality iff $a$
is non-increasing.
\end{lemma}
\begin{proof}
$\Phi(X_a)=\sum_j\sum_{i=1}^{a_j}(i+j)=\sum_j\binom{a_j+1}{2}+F(a)$ with $F(a)=\sum_j j\,a_j$.
The first sum does not change when $a$ is permuted.  If $a_j<a_{j+1}$, exchanging these two entries
changes $F$ by $j a_{j+1}+(j+1)a_j-ja_j-(j+1)a_{j+1}=a_j-a_{j+1}<0$.  Starting from $a$, repeatedly exchange
some adjacent pair with $a_j<a_{j+1}$.  Each exchange lowers the non-negative integer $F$, so the process
stops, and it stops at a sequence with no adjacent ascent, i.e.\ at the non-increasing rearrangement
$a^*$ (which is unique).  If $a$ is not non-increasing at least one exchange happens, so $F(a^*)<F(a)$; if
$a$ is non-increasing then $a=a^*$.
\end{proof}

\begin{corollary}\label{cor:energy}
For $\lambda\vdash n$: $\Phi(Y(\B(\lambda)))\le\Phi(Y(\lambda))$, with equality iff $\ell(\lambda)\ge\lambda_1-1$;
in the case of equality $Y(\B(\lambda))=\rho(Y(\lambda))$.
\end{corollary}
\begin{proof}
By Lemma~\ref{lem:rho}, $\rho$ is injective and preserves $i+j$, so $\Phi(\rho(X))=\Phi(X)$.  With $a$ as in
Lemma~\ref{lem:move}: $\Phi(Y(\lambda))=\Phi(X_a)\ge\Phi(X_{a^*})=\Phi(Y(\B(\lambda)))$, with equality iff $a$ is
non-increasing (Lemma~\ref{lem:rearr}), i.e.\ iff $\ell(\lambda)\ge\lambda_1-1$, and then
Lemma~\ref{lem:move} gives $Y(\B(\lambda))=\rho(Y(\lambda))$.
\end{proof}

For $k\ge1$ and $E\subseteq D_k$ let $\lambda_E$ be the partition whose $j$-th part, $1\le j\le k$, is
$k-j+[d_k(j)\in E]$ (zero parts omitted).  The sequence $(k-j+\varepsilon_j)_{j=1}^k$ with
$\varepsilon_j\in\{0,1\}$ is non-increasing (consecutive differences are $1+\varepsilon_j-\varepsilon_{j+1}\ge0$),
so $\lambda_E$ is a partition, $|\lambda_E|=T_{k-1}+|E|$, and $Y(\lambda_E)=\Delta_{k-1}\cup E$ (column $j$ consists
of the $k-j$ cells of $\Delta_{k-1}$ and, if $d_k(j)\in E$, the cell $d_k(j)=(k+1-j,j)$).  Writing
$e_{k-j}=[d_k(j)\in E]$ we get $\lambda_E=(k-1+e_{k-1},\,k-2+e_{k-2},\dots,1+e_1,\,e_0)$.

\begin{theorem}[Brandt \cite{Br}; see also \cite{AD,Et,GH}]\label{thm:brandt}
Let $n=T_{k-1}+r$ with $1\le r\le k$.  A partition of $n$ is cyclic if and only if it equals $\lambda_E$
for some $E\subseteq D_k$ with $|E|=r$; equivalently, iff it has the form
$(k-1+e_{k-1},\dots,1+e_1,e_0)$ with $e_i\in\{0,1\}$ and $\sum_ie_i=r$ (zero parts omitted).
Moreover $Y(\B(\lambda_E))=\rho(Y(\lambda_E))$, i.e.\ $\B(\lambda_E)=\lambda_{\rho(E)}$, and $\B^k(\lambda_E)=\lambda_E$.
\end{theorem}
\begin{proof}
($\Leftarrow$ and ``moreover'').  Let $\mu=\lambda_E$ with $E\subseteq D_k$, $|E|=r$.  Row~1 of
$Y(\mu)=\Delta_{k-1}\cup E$ consists of the cells $(1,w)=d_w(w)$ with $w\le k-1$ and possibly $d_k(k)$;
so $\ell(\mu)=k-1+[d_k(k)\in E]\ge k-1$.  Column~1 consists of $d_w(1)$ for $w\le k-1$ and possibly $d_k(1)$;
so $\mu_1=k-1+[d_k(1)\in E]\le k$.  Thus $\ell(\mu)\ge\mu_1-1$ and Corollary~\ref{cor:energy} and
Lemma~\ref{lem:rho} give $Y(\B(\mu))=\rho(\Delta_{k-1})\cup\rho(E)=\Delta_{k-1}\cup\rho(E)$ with
$\rho(E)\subseteq D_k$, $|\rho(E)|=r$.  Hence $\B(\lambda_E)=\lambda_{\rho(E)}$, and by induction
$\B^k(\lambda_E)=\lambda_{\rho^k(E)}=\lambda_E$ because $\rho^k$ is the identity on $D_k$.  So $\mu$ is cyclic.

($\Rightarrow$).  Let $\lambda\vdash n$ with $\B^P(\lambda)=\lambda$, $P\ge 1$, and write
$\lambda^{(m)}=\B^m(\lambda)$.  By Corollary~\ref{cor:energy},
$\Phi(Y(\lambda^{(0)}))\ge\Phi(Y(\lambda^{(1)}))\ge\dots\ge\Phi(Y(\lambda^{(P)}))=\Phi(Y(\lambda^{(0)}))$, so all these are
equalities and $Y(\lambda^{(m+1)})=\rho(Y(\lambda^{(m)}))$ for $0\le m<P$.  As $\lambda^{(m+P)}=\lambda^{(m)}$, this holds for
all $m\ge0$, and therefore $Y(\lambda^{(m)})=\rho^m(Y(\lambda))$ for all $m\ge 0$.

\emph{Claim: there is no $w\ge1$ such that $D_w\not\subseteq Y(\lambda)$ and $D_{w+1}\cap Y(\lambda)\ne\emptyset$.}
Suppose $d_w(\alpha)\notin Y(\lambda)$ and $d_{w+1}(\beta)\in Y(\lambda)$.  Since $\gcd(w,w+1)=1$, the Chinese remainder
theorem gives $m\ge0$ with $m\equiv w-\alpha\pmod w$ and $m\equiv w+1-\beta\pmod{w+1}$.  By Lemma~\ref{lem:rho},
$\rho^m(d_w(\alpha))=d_w(w)=(1,w)$ and $\rho^m(d_{w+1}(\beta))=d_{w+1}(w+1)=(1,w+1)$.  As $\rho^m$ is injective,
$(1,w)\notin\rho^m(Y(\lambda))=Y(\lambda^{(m)})$ while $(1,w+1)\in Y(\lambda^{(m)})$, contradicting \eqref{eq:Y1}.

$Y(\lambda)$ is finite, so some $D_w$ is not contained in it; let $K\ge1$ be the least such $w$.  Then
$\Delta_{K-1}\subseteq Y(\lambda)$.  By the claim, $D_{K+1}\cap Y(\lambda)=\emptyset$; then $D_{K+1}\not\subseteq Y(\lambda)$, so
again by the claim $D_{K+2}\cap Y(\lambda)=\emptyset$, and inductively $D_w\cap Y(\lambda)=\emptyset$ for all $w>K$.  Thus
$Y(\lambda)=\Delta_{K-1}\cup E'$ with $E'=Y(\lambda)\cap D_K\ne D_K$; let $r'=|E'|\in\{0,\dots,K-1\}$, so
$n=T_{K-1}+r'$ and $Y(\lambda)=Y(\lambda_{E'})$ (with $k$ replaced by $K$ in the definition of $\lambda_{E'}$).
If $r'\ge1$, then $T_{K-1}<n<T_K$, so by Lemma~\ref{lem:rank} $K=k$, $r'=r$, and $\lambda=\lambda_{E'}$.
If $r'=0$, then $K\ge2$ (as $n\ge1$) and $n=T_{K-1}=T_{K-2}+(K-1)$, so $k=K-1$, $r=k$, and
$Y(\lambda)=\Delta_k=\Delta_{k-1}\cup D_k$, i.e.\ $\lambda=\lambda_{D_k}$ with $|D_k|=k=r$.
\end{proof}

\begin{corollary}\label{cor:cyc}
Let $n=T_{k-1}+r$, $1\le r\le k$.
\begin{enumerate}
\item[(a)] If $r=k$ (i.e.\ $n=T_k$), the only cyclic partition of $n$ is $\delta_k=(k,k-1,\dots,1)$, and
$\B(\delta_k)=\delta_k$.
\item[(b)] Every cyclic $\mu\vdash n$ satisfies $Y(\mu)\subseteq\Delta_k$, $\Delta_{k-1}\subseteq Y(\mu)$,
$\ell(\mu)\in\{k-1,k\}$ and $\B^k(\mu)=\mu$.
\item[(c)] If $\mu\vdash n$ is cyclic, then for every $a\ge0$:
$\sum_{i=a}^{a+k-1}\ell(\B^i(\mu))=k(k-1)+r$.  If moreover $r<k$, both values $k-1$ and $k$ occur among
$\ell(\B^a(\mu)),\dots,\ell(\B^{a+k-1}(\mu))$.
\end{enumerate}
\end{corollary}
\begin{proof}
(a) and (b) are immediate from Theorem~\ref{thm:brandt} ($E=D_k$ when $r=k$; $\ell(\lambda_E)=k-1+[d_k(k)\in E]$ was
computed in its proof).  (c) Write $\mu=\lambda_E$.  Then $\B^i(\mu)=\lambda_{\rho^i(E)}$, so
$\ell(\B^i(\mu))=k-1+[d_k(k)\in\rho^i(E)]=k-1+[d_k(j_i)\in E]$, where $d_k(j_i)=\rho^{-i}(d_k(k))$, i.e.\
$j_i\equiv k-i\pmod k$ (Lemma~\ref{lem:rho}).  As $i$ runs through $a,\dots,a+k-1$, $j_i$ runs through all
of $1,\dots,k$ once, so the sum is $k(k-1)+|E|=k(k-1)+r$.  If $0<r<k$, some $j_i$ lies in $E$ and some does not.
\end{proof}

\section{Piles and the part-count sequence}\label{sec:piles}

This is the ``$\mathrm{diagram}_B(\lambda)$'' bookkeeping of Griggs--Ho \cite[\S3]{GH}, written with
explicitly labelled piles.

Fix $\lambda\vdash n$, write $\lambda^{(i)}=\B^i(\lambda)$ and
\[c_i=\ell(\lambda^{(i-1)})\qquad(i\ge1).\]
We call $c(\lambda)=(c_1,c_2,\dots)$ the \emph{part-count sequence}; clearly
$c(\lambda^{(a)})=(c_{a+1},c_{a+2},\dots)$.  Every $c_i\ge1$.

\begin{lemma}[piles]\label{lem:piles}
For every $i\ge0$ the multiset of parts of $\lambda^{(i)}$ is
\[\{\lambda_m-i:\ 1\le m\le\ell(\lambda),\ \lambda_m>i\}\ \uplus\ \{c_j-(i-j):\ 1\le j\le i,\ c_j>i-j\}.\]
\end{lemma}
\begin{proof}
Induction on $i$; $i=0$ is clear.  $\lambda^{(i+1)}=\B(\lambda^{(i)})$ has the part $\ell(\lambda^{(i)})=c_{i+1}=c_{i+1}-0$ and
the parts of $\lambda^{(i)}$ decreased by one, the zeros being discarded: this is the stated multiset for
$i+1$ (the term $j=i+1$ is $c_{i+1}$).
\end{proof}

We use \emph{labels}: $O_1,\dots,O_{\ell(\lambda)}$ (the original piles) and $N_1,N_2,\dots$ ($N_j$ is the pile
created by the $j$-th move, of size $c_j$ when created).  By Lemma~\ref{lem:piles}, at time $i$ (i.e.\ in
$\lambda^{(i)}$) the pile $O_m$ has size $\lambda_m-i$ and the pile $N_j$ ($j\le i$) has size $c_j-(i-j)$, a pile
being present iff its size is positive.  For $i\ge1$ let $\R_i$ be the set of labels of the piles present
at time $i-1$:
\begin{equation}\label{eq:alive}
O_m\in \R_i\iff\lambda_m\ge i,\qquad N_j\in \R_i\iff j\le i-1\ \text{and}\ c_j\ge i-j .
\end{equation}
By Lemma~\ref{lem:piles} (with $i-1$), $|\R_i|=\ell(\lambda^{(i-1)})=c_i$.  Let $Z_i=\R_i\setminus \R_{i+1}$ (the piles
of size $1$ at time $i-1$, which disappear in move $i$).

\begin{lemma}[basic facts]\label{lem:basic}
For all $i\ge1$:
\begin{enumerate}
\item[(a)] $N_i\in \R_{i+1}$ and $N_i\notin \R_i$.
\item[(b)] $\R_{i+1}=(\R_i\setminus Z_i)\cup\{N_i\}$, hence $c_{i+1}=c_i+1-|Z_i|$ and $c_{i+1}\le c_i+1$.
\item[(c)] If $c_{i+1}=c_i+1$ then $Z_i=\emptyset$, i.e.\ $\R_i\subseteq \R_{i+1}$.  If $c_{i+1}=c_i$ then $|Z_i|=1$.
\item[(d)] For $j<i$: $N_j\in \R_i$ implies $N_j\in \R_{i'}$ for $j<i'\le i$, and $N_j\notin \R_i$ implies
$N_j\notin \R_{i'}$ for all $i'\ge i$.  Likewise for the labels $O_m$.
\end{enumerate}
\end{lemma}
\begin{proof}
(a) $c_i\ge1=(i+1)-i$; and labels $N_j$ in $\R_i$ have $j\le i-1$.  (d) For $i'>j$ the clause $j\le i'-1$ in
\eqref{eq:alive} holds automatically, so membership of $N_j$ in $\R_{i'}$ is the condition $c_j\ge i'-j$, which
is monotone in $i'$; for $O_m$ the condition $\lambda_m\ge i'$ is monotone as well.  (b) By (d), $\R_{i+1}\setminus\{N_i\}\subseteq \R_i$, so
$\R_{i+1}=(\R_i\cap \R_{i+1})\cup\{N_i\}=(\R_i\setminus Z_i)\cup\{N_i\}$, a disjoint union by (a); take cardinalities.
(c) follows from $c_{i+1}=c_i+1-|Z_i|$.
\end{proof}

\begin{lemma}[gap lemma]\label{lem:gap}
Let $1\le j<i$.  If $N_j\notin \R_i$ and $N_{j+1}\in \R_{i+1}$, then $c_j=i-j-1$, $c_{j+1}=i-j$, and
$N_{j+1}\notin \R_{i+2}$.
\end{lemma}
\begin{proof}
By \eqref{eq:alive}, $c_j\le i-j-1$ and $c_{j+1}\ge (i+1)-(j+1)=i-j$.  By Lemma~\ref{lem:basic}(b),
$c_{j+1}\le c_j+1\le i-j$.  Hence both equalities hold, and $c_{j+1}=i-j<(i+2)-(j+1)$ gives
$N_{j+1}\notin \R_{i+2}$.
\end{proof}

\begin{lemma}[the sequence determines the partition]\label{lem:determ}
For $i\ge1$, $\#\{m:\lambda_m\ge i\}=c_i-\#\{j:1\le j\le i-1,\ c_j\ge i-j\}$.  Consequently, if two partitions
have the same part-count sequence, they are equal.
\end{lemma}
\begin{proof}
The formula is $|\R_i|=c_i$ together with \eqref{eq:alive}.  The numbers $\#\{m:\lambda_m\ge i\}$, $i\ge1$,
determine $\lambda$, since $\lambda_m=\#\{i\ge1:\#\{m':\lambda_{m'}\ge i\}\ge m\}$.
\end{proof}

\begin{lemma}[window sums; generalising {\cite[Prop.~3.2(2)]{GH}}]\label{lem:window}
For all $a\ge0$ and $h\ge1$: $c_{a+1}+c_{a+2}+\dots+c_{a+h}\le n+T_{h-1}$.
\end{lemma}
\begin{proof}
Let $\mu=\lambda^{(a)}$; its part-count sequence is $(c_{a+1},c_{a+2},\dots)$.  Applying \eqref{eq:alive} to $\mu$,
for $1\le g\le h$ the number $c_{a+g}$ equals $\#\{m:\mu_m\ge g\}$ plus the number of piles created in the
first $g-1$ moves starting from $\mu$ that are still present, which is at most $g-1$.  Summing,
$\sum_{g=1}^h c_{a+g}\le\sum_{g\ge1}\#\{m:\mu_m\ge g\}+\sum_{g=1}^h(g-1)=|\mu|+T_{h-1}=n+T_{h-1}$.
\end{proof}

\section{Ramps}\label{sec:ramps}

The following lemmas are Lemmas 3.4--3.6 and Proposition 3.2(3) of Griggs--Ho \cite{GH}; we give complete
proofs (the proof of \cite[Lemma~3.5]{GH} ends with ``continue this process'', which we replace by an
explicit induction) and we add the equality statement in Proposition~\ref{prop:rampbound}.

\begin{definition}
A \emph{ramp} of $\lambda$ is a triple $(p,q,x)$ of integers with $p\ge1$, $q\ge p+2$ and
\[c_p=x-1,\qquad c_{p+1}=c_{p+2}=\dots=c_{q-1}=x,\qquad c_q=x+1 .\]
Its \emph{length} is $q-p$ and its \emph{value} is $x$ (note $x\ge2$, since $c_p\ge1$).
\end{definition}

\begin{lemma}[sandwich]\label{lem:sandwich}
If $1\le i<j$ and $c_i<x<c_j$ for an integer $x$, then there is a ramp $(p,q,x)$ with $i\le p<q\le j$.
\end{lemma}
\begin{proof}
Let $p$ be the largest $m$ with $i\le m<j$ and $c_m\le x-1$ (it exists: $m=i$ qualifies).  By
Lemma~\ref{lem:basic}(b), $c_{p+1}\le c_p+1\le x$.  Since $c_j>x$, this rules out $p+1=j$; so $p+1<j$, and
$c_{p+1}\ge x$ by maximality of $p$.  Hence $c_p=x-1$ and $c_{p+1}=x$.  Let $q$ be the least $m$ with $p<m\le j$ and
$c_m\ge x+1$ (it exists: $m=j$ qualifies).  Then $q\ge p+2$ since $c_{p+1}=x$.  For $p<m<q$ we have $c_m\le x$
(minimality of $q$) and $c_m\ge x$ (maximality of $p$, as $m<j$), so $c_m=x$.  Finally
$x+1\le c_q\le c_{q-1}+1=x+1$.
\end{proof}

\begin{lemma}[short ramps {\cite[Lemma~3.6]{GH}}]\label{lem:short}
If $(p,p+2,x)$ is a ramp, then $p\le x$.
\end{lemma}
\begin{proof}
Suppose $p\ge x+1$; then $p\ge3$.  Since $|\R_p|=c_p=x-1\le p-2$, not all of $N_1,\dots,N_{p-1}$ lie in
$\R_p$, while $N_{p-1}\in \R_p$ (Lemma~\ref{lem:basic}(a)).  Let $j$ be the largest index in $[1,p-2]$ with
$N_j\notin \R_p$; then $N_{j+1}\in \R_p$.  As $c_{p+1}=c_p+1$ and $c_{p+2}=c_{p+1}+1$, Lemma~\ref{lem:basic}(c)
gives $\R_p\subseteq \R_{p+1}\subseteq \R_{p+2}$, so $N_{j+1}\in \R_{p+1}$ and $N_{j+1}\in \R_{p+2}$.  But
Lemma~\ref{lem:gap} (with $i=p$) gives $N_{j+1}\notin \R_{p+2}$, a contradiction.
\end{proof}

\begin{lemma}[descent {\cite[Lemma~3.5]{GH}}]\label{lem:descent}
Let $(p,q,x)$ be a ramp with $q-p\ge3$ and $p\ge x+1$.  Then there is a ramp $(p',q',x')$ with
\[p-x\le p'\le p-2,\qquad x'=p-p'\le x,\qquad 2\le q'-p'\le q-p-1 .\]
\end{lemma}
\begin{proof}
The $x$ labels $N_{p-x},\dots,N_{p-1}$ exist since $p-x\ge1$.  As $|\R_p|=x-1$, one of them is not in $\R_p$;
$N_{p-1}\in \R_p$.  Let $p'$ be the largest index in $[p-x,p-2]$ with $N_{p'}\notin \R_p$ and put $y=p-p'$, so
$2\le y\le x$ and $N_{p'+1},\dots,N_{p-1}\in \R_p$.  As $c_{p+1}=c_p+1$, $\R_p\subseteq \R_{p+1}$, and with
Lemma~\ref{lem:basic}(a), $N_{p'+1},\dots,N_p\in \R_{p+1}$.  Lemma~\ref{lem:gap} ($j=p'$, $i=p$) gives
$c_{p'}=y-1$ and $c_{p'+1}=y$.

For $1\le l\le q-p-1$ let $H(l)$ be the statement
\[H(l):\qquad c_{p'+1}=\dots=c_{p'+l}=y\quad\text{and}\quad N_{p'+l},N_{p'+l+1},\dots,N_{p+l-1}\in \R_{p+l}.\]
We just proved $H(1)$.  Assume $H(l)$.  Since
$c_{p'+l}=y=p-p'$, \eqref{eq:alive} gives, for $i>p'+l$: $N_{p'+l}\in \R_i\iff i\le p+l$.  Hence
$N_{p'+l}\in Z_{p+l}$.
\begin{itemize}
\item If $l=q-p-1$, then $p+l=q-1$ and $c_q=c_{q-1}+1$, so $Z_{q-1}=\emptyset$ by Lemma~\ref{lem:basic}(c): a
contradiction.  So $H(q-p-1)$ is false.
\item If $l\le q-p-2$, then $p+l$ and $p+l+1$ lie in $[p+1,q-1]$, so $c_{p+l+1}=c_{p+l}=x$ and $|Z_{p+l}|=1$,
i.e.\ $Z_{p+l}=\{N_{p'+l}\}$.  By Lemma~\ref{lem:basic}(b),
$\R_{p+l+1}=(\R_{p+l}\setminus\{N_{p'+l}\})\cup\{N_{p+l}\}\supseteq\{N_{p'+l+1},\dots,N_{p+l}\}$.  In particular
$N_{p'+l+1}\in \R_{p+l+1}$, i.e.\ $c_{p'+l+1}\ge(p+l+1)-(p'+l+1)=y$, and $c_{p'+l+1}\le c_{p'+l}+1=y+1$.  So either
$c_{p'+l+1}=y+1$, or $c_{p'+l+1}=y$ and then $H(l+1)$ holds.
\end{itemize}
Let $l_0$ be the largest $l\le q-p-1$ such that $H(1),\dots,H(l)$ all hold; then $1\le l_0\le q-p-2$, and
since $H(l_0+1)$ fails, the second item gives $c_{p'+l_0+1}=y+1$.  Thus $(p',p'+l_0+1,y)$ is a ramp of
length $l_0+1\in[2,q-p-1]$, and $x'=y$ has the stated properties.
\end{proof}

\begin{proposition}[ramp bound]\label{prop:rampbound}
Every ramp $(p,q,x)$ satisfies $p\le(q-p-1)\,x$.  If equality holds, then $(jx,\ jx+j+1,\ x)$ is a ramp
for every $j=1,2,\dots,q-p-1$.
\end{proposition}
\begin{proof}
Strong induction on the length $L=q-p\ge2$.  If $L=2$, Lemma~\ref{lem:short} gives $p\le x$; if $p=x$, the
ramp itself is $(x,x+2,x)$.  Let $L\ge3$.  If $p\le x$ then $p\le x<(L-1)x$.  Otherwise
Lemma~\ref{lem:descent} provides a ramp $(p',q',x')$ with $p'\ge p-x$, $x'\le x$ and length $L'\le L-1$; by
induction $p'\le(L'-1)x'$, hence
\[p\le p'+x\le(L'-1)x'+x\le(L-2)x+x=(L-1)x .\]
If $p=(L-1)x$, all these inequalities are equalities: $p'=p-x=(L-2)x$ and $(L'-1)x'=(L-2)x$; since
$1\le L'-1\le L-2$ and $x'\le x$, this forces $L'=L-1$ and $x'=x$.  So $((L-2)x,(L-2)x+L-1,x)$ is a ramp
attaining equality, and by induction $(jx,jx+j+1,x)$ is a ramp for $j=1,\dots,L-2$; for $j=L-1$ it is
$(p,q,x)$ itself.
\end{proof}

\begin{lemma}[long ramps {\cite[Lemma~3.4]{GH}}]\label{lem:long}
If $(p,q,x)$ is a ramp with $q-p=x+1$, then $q\le n+1$.  (Recall that $x\ge2$ for every ramp.)
\end{lemma}
\begin{proof}
For $p+1\le j\le q-1$ we have $c_j=x\ge q-j$, so $N_j\in \R_q$: these are $x$ piles.  Also
$c_p=x-1<q-p$, so $N_p\notin \R_q$.  Since $|\R_q|=c_q=x+1$, $\R_q$ contains exactly one further label $P$,
and $P=O_m$ for some $m$, or $P=N_j$ with $1\le j\le p-1$.
If $P=O_m$, then $q\le\lambda_m\le n$.  If $P=N_1$, then $c_1\ge q-1$, and $c_1=\ell(\lambda)\le n$.
Suppose $P=N_j$ with $2\le j\le p-1$.  Then $N_{j-1}\notin \R_q$.  As $c_q=c_{q-1}+1$, $\R_{q-1}\subseteq \R_q$, so
$N_{j-1}\notin \R_{q-1}$.  Next, $q-2\ge p+1$ (as $x\ge2$), so $c_{q-1}=c_{q-2}=x$ and $|Z_{q-2}|=1$; since
$c_p=x-1=(q-2)-p$, $N_p\in \R_{q-2}\setminus \R_{q-1}$, so $Z_{q-2}=\{N_p\}$.  As $j-1\ne p$, $N_{j-1}\in \R_{q-2}$ would give
$N_{j-1}\in \R_{q-1}$ (Lemma~\ref{lem:basic}(b)); hence $N_{j-1}\notin \R_{q-2}$, i.e.\ $c_{j-1}\le q-j-2$.  But
$N_j\in \R_q$ gives $c_j\ge q-j$, so $c_j\ge c_{j-1}+2$, contradicting Lemma~\ref{lem:basic}(b).
\end{proof}

\section{The tail lemma}\label{sec:tail}

This is \cite[Lemma~3.3(2) and Lemma~4.3(2)]{GH}, in a form that serves both the triangular and the
non-triangular case.

\begin{lemma}\label{lem:tail}
Let $k\ge2$ and $t\ge k+1$ with $c_t=k-1$, $c_{t+1}=k$, and $c_j\ge k-1$ for $t\le j\le t+k-2$.  Then at
least one of the following holds.
\begin{enumerate}
\item[(A)] There is a ramp $(p,q,k)$ with $t-k\le p<q\le t-1$.
\item[(B)] There is a ramp $(p,q,k-1)$ with $t-k+1\le p<q\le t+1$.
\item[(C)] Every part of $\lambda^{(t-1)}$ is at most $k$, and $c_{t+m}=c_{t-k+m}$ for $0\le m\le k-1$.
\end{enumerate}
\end{lemma}
\begin{proof}
The $k$ labels $N_{t-k},\dots,N_{t-1}$ exist ($t-k\ge1$); as $|\R_t|=k-1$, one of them is not in $\R_t$, and
$N_{t-1}\in \R_t$.  Let $i$ be the largest index in $[t-k,t-2]$ with $N_i\notin \R_t$.  Then
$N_{i+1},\dots,N_{t-1}\in \R_t$; as $c_{t+1}=c_t+1$, $\R_t\subseteq \R_{t+1}$, so $N_{i+1}\in \R_{t+1}$, and
Lemma~\ref{lem:gap} gives $c_i=t-i-1$ and $c_{i+1}=t-i$.

\emph{Case 1: $i\ge t-k+1$.}  Then $c_i\le k-2<k-1<k=c_{t+1}$, and Lemma~\ref{lem:sandwich} gives a ramp
$(p,q,k-1)$ with $i\le p<q\le t+1$: (B) holds.

\emph{Case 2: $i=t-k$.}  Then $c_{t-k}=k-1$, $c_{t-k+1}=k$, and $\R_t=\{N_{t-k+1},\dots,N_{t-1}\}$ (these $k-1$ labels lie in
$\R_t$ and $|\R_t|=k-1$).  If $c_j\le k-2$ for some $j\in[t-k+2,t-1]$, Lemma~\ref{lem:sandwich} applied to $j<t+1$ gives
(B).  If $c_j\ge k+1$ for some $j\in[t-k+2,t-1]$, Lemma~\ref{lem:sandwich} applied to $t-k<j$ (with
$c_{t-k}=k-1<k<c_j$) gives a ramp $(p,q,k)$ with $t-k\le p<q\le j\le t-1$: (A) holds.  Otherwise
$c_j\in\{k-1,k\}$ for all $j\in[t-k+1,t-1]$; we show (C).  The parts of $\lambda^{(t-1)}$ are the sizes
$c_j-(t-1-j)\le k$ of the piles $N_j\in \R_t$.  Write $c_{t-1-u}=k-1+\varepsilon_u$ with $\varepsilon_u\in\{0,1\}$
($0\le u\le k-2$).  Fix $m\in[0,k-1]$ and determine $\R_{t+m}$:
\begin{itemize}
\item labels not in $\R_t$ (all $O$'s and all $N_j$ with $j\le t-k$) are not in $\R_{t+m}$
(Lemma~\ref{lem:basic}(d));
\item for $t\le j\le t+m-1$: $c_j\ge k-1\ge m\ge t+m-j$, so $N_j\in \R_{t+m}$ ($m$ labels);
\item $N_{t-1-u}\in \R_{t+m}\iff k-1+\varepsilon_u\ge m+1+u\iff u\le k-2-m+\varepsilon_u$.
\end{itemize}
For $m=0$ this gives $c_t=k-1=c_{t-k}$.  For $1\le m\le k-1$ the labels $N_{t-1-u}$ with $u\le k-2-m$ are
in $\R_{t+m}$, the one with $u=k-1-m$ is in iff $\varepsilon_{k-1-m}=1$, and those with $u\ge k-m$ are not.  So
$c_{t+m}=m+(k-1-m)+\varepsilon_{k-1-m}=k-1+\varepsilon_{k-1-m}=c_{t-k+m}$.
\end{proof}

\section{Two lemmas for non-triangular $n$}\label{sec:c3lemmas}

\begin{lemma}[too many parts]\label{lem:overfull}
Let $n=T_{k-1}+r$ with $1\le r\le k$ and $\lambda\vdash n$.  If $c_m\ge k+1$ for some $m\ge1$, then
$\dB(\lambda)\ge m+1$.
\end{lemma}
\begin{proof}
$\lambda^{(m-1)}$ has $c_m\ge k+1$ parts, so it is not cyclic (Corollary~\ref{cor:cyc}(b)), and $\dB(\lambda)\ge m$.
Suppose $\dB(\lambda)=m$, i.e.\ $\lambda^{(m)}$ is cyclic.  Its iterates are $\lambda^{(m)},\dots,\lambda^{(m+k-1)}$, whose part
counts are $c_{m+1},\dots,c_{m+k}$, so Corollary~\ref{cor:cyc}(c) gives $\sum_{i=m+1}^{m+k}c_i=k(k-1)+r$.
Lemma~\ref{lem:window} with $a=m-1$, $h=k$ gives $\sum_{i=m}^{m+k-1}c_i\le n+T_{k-1}=k(k-1)+r$.  Subtracting,
$c_m\le c_{m+k}\le k$, a contradiction.
\end{proof}

\begin{lemma}[the anchor $t$]\label{lem:anchor}
Let $k\ge2$, $n=T_{k-1}+r$ with $1\le r\le k-1$, and $\lambda\vdash n$.  The set
\[\mathcal T=\{i\ge1:\ \lambda^{(i-1)}\text{ is cyclic},\ c_i=k-1,\ c_{i+1}=k\}\]
is non-empty; let $t=\min\mathcal T$.  Then $\dB(\lambda)\le t-1$, and for all $j\ge t$:
$c_j\in\{k-1,k\}$, $c_{j+k}=c_j$, and $\sum_{i=j}^{j+k-1}c_i=k(k-1)+r$.
\end{lemma}
\begin{proof}
Let $d=\dB(\lambda)$.  For $i\ge d+1$, $\lambda^{(i-1)}$ is cyclic, so by Corollary~\ref{cor:cyc}(b),(c) the
sequence $(c_i)_{i\ge d+1}$ takes only the values $k-1$ and $k$, is periodic with period $k$, and takes both
values.  Choose $i_0\ge d+1$ with $c_{i_0}=k-1$, and let $i_1>i_0$ be least with $c_{i_1}=k$.  Then
$c_{i_1-1}=k-1$, and $i=i_1-1\ge d+1$ lies in $\mathcal T$.  If $t\in\mathcal T$, then $\lambda^{(t-1)}$ is cyclic, so
$d\le t-1$.  The remaining statements hold because every $\lambda^{(j-1)}$ with $j\ge t$ is cyclic, by
Corollary~\ref{cor:cyc}(b),(c).
\end{proof}

\section{The upper bound}\label{sec:c3upper}

\begin{theorem}[{\cite[Thm.~4.4(1)]{GH}}]\label{thm:c3upper}
Let $k\ge4$ and $n=T_{k-1}+r$ with $1\le r\le k-1$.  Then $\DB(n)\le k^2-2k-1$.
\end{theorem}
\begin{proof}
\emph{$k=4$} ($n\in\{7,8,9\}$, bound $7$): by exhaustive computation over the $15+22+30$ partitions,
$\DB(7)=4$, $\DB(8)=5$, $\DB(9)=7$ (Appendix~\ref{app:code}; this matches \cite[Fig.~1]{GH}).
The general argument below does not cover $k=4$: in case (B) with $L=k$ it only gives $d\le T_k-2=8>7$.

\emph{$k\ge5$.}  Let $\lambda\vdash n$, $d=\dB(\lambda)$, and let $t$ be as in Lemma~\ref{lem:anchor}, so $d\le t-1$.  If
$t\le k$, then $d\le k-1<k^2-2k-1$.  Assume $t\ge k+1$.  By Lemma~\ref{lem:anchor} the hypotheses of
Lemma~\ref{lem:tail} hold.

\emph{(C) is impossible.}  Suppose $c_{t-k+m}=c_{t+m}$ for $0\le m\le k-1$.  For $m\ge k$,
$c_{t-k+m}=c_{t+(m-k)}=c_{t+m}$ by the periodicity in Lemma~\ref{lem:anchor}.  Thus $\lambda^{(t-k-1)}$ and
$\lambda^{(t-1)}$ have the same part-count sequence $(c_t,c_{t+1},\dots)$, so they are equal
(Lemma~\ref{lem:determ}).  In particular $\lambda^{(t-k-1)}$ is cyclic.  Since also $c_{t-k}=c_t=k-1$,
$c_{t-k+1}=c_{t+1}=k$ and $t-k\ge1$, we get $t-k\in\mathcal T$, contradicting the minimality of $t$.

\emph{(A): a ramp $(p,q,k)$ with $t-k\le p<q\le t-1$}, of length $L=q-p\le k-1$.
If $L=k-1$, then $p=t-k$, $q=t-1$ and $c_{t-1}=k+1$.  Lemma~\ref{lem:window} ($a=t-2$, $h=k$) gives
$c_{t-1}+\sum_{i=t}^{t+k-2}c_i\le k(k-1)+r$, while $\sum_{i=t}^{t+k-1}c_i=k(k-1)+r$ by Lemma~\ref{lem:anchor}.
Hence $c_{t-1}\le c_{t+k-1}\le k$, a contradiction.  If $L\le k-2$, Proposition~\ref{prop:rampbound} gives
$p\le(k-3)k$, and $d\le t-1\le p+k-1\le k^2-2k-1$.

\emph{(B): a ramp $(p,q,k-1)$ with $t-k+1\le p<q\le t+1$}, of length $L\le k$.
\begin{itemize}
\item $L=k$: then $(p,q)=(t-k+1,t+1)$, and Lemma~\ref{lem:long} (with $x=k-1\ge2$) gives $t+1\le n+1$.  So
$d\le t-1\le n-1\le T_{k-1}+k-2=T_k-2$, and $(k^2-2k-1)-(T_k-2)=\tfrac12(k^2-5k+2)>0$ for $k\ge5$.
\item $L=k-1$: then $(p,q)=(t-k+1,t)$ or $(t-k+2,t+1)$.  The first would give $c_t=k$, but $c_t=k-1$.  In the
second, $c_{t-k+2}=k-2=t-(t-k+2)$, so by \eqref{eq:alive} $N_{t-k+2}\in \R_t\setminus \R_{t+1}=Z_t$.  But
$c_{t+1}=c_t+1$ forces $Z_t=\emptyset$ (Lemma~\ref{lem:basic}(c)).  Contradiction.
\item $L\le k-2$: $p\le(k-3)(k-1)$ by Proposition~\ref{prop:rampbound}, and
$d\le t-1\le p+k-2\le k^2-3k+1<k^2-2k-1$.
\end{itemize}
In all possible cases $d\le k^2-2k-1$.
\end{proof}

\section{Equality for $n=T_k-1$}\label{sec:c3lower}

For $k\ge3$, $h\in\{1,\dots,k\}$ and $x\in\{1,\dots,k+1\}$ put
\[X'(h,x)=\Delta_{k-1}\cup\bigl(D_k\setminus\{d_k(h),d_k(h\oplus_k1)\}\bigr)\cup\{d_{k+1}(x)\},\]
where for $w\ge1$ and $1\le j\le w$, $j\oplus_w1=j+1$ if $j<w$ and $w\oplus_w1=1$.

\begin{lemma}\label{lem:c3conf}
Let $k\ge3$ and let $\mu$ be a partition with $Y(\mu)=X'(h,x)$.  If $(h,x)\ne(k-1,1)$, then
$Y(\B(\mu))=X'(h\oplus_k1,\,x\oplus_{k+1}1)$.  If $(h,x)=(k-1,1)$, then $\mu=(k+1,k-1,k-2,\dots,3,1)$ (to be read as $(4,1)$ when $k=3$) and
$\B(\mu)=(k,k-1,\dots,2)$, which is cyclic.
\end{lemma}
\begin{proof}
Let $H=\{h,h\oplus_k1\}$, a set of two elements as $k\ge3$.  The cells of row~1 are among the
$d_w(w)=(1,w)$ and those of column~1 among the $d_w(1)=(w,1)$.  Counting them,
$\ell(\mu)=k-1+[k\notin H]+[x=k+1]$ and $\mu_1=k-1+[1\notin H]+[x=1]$.  So
$\ell(\mu)\le\mu_1-2$ iff $k\in H$, $x\ne k+1$, $1\notin H$ and $x=1$.  Now $k\in H$ iff $h\in\{k-1,k\}$, and
$h=k$ gives $H=\{k,1\}\ni1$.  So this happens iff $(h,x)=(k-1,1)$ (note $1\notin\{k-1,k\}$ as $k\ge3$).  If
$(h,x)\ne(k-1,1)$, Corollary~\ref{cor:energy} and Lemma~\ref{lem:rho} give $Y(\B(\mu))=\rho(X'(h,x))=
X'(h\oplus_k1,x\oplus_{k+1}1)$.  If $(h,x)=(k-1,1)$, the column lengths are $k+1$ (column 1), $k+1-j$ for
$2\le j\le k-2$, $1$ (column $k-1$: the cell of $\Delta_{k-1}$ only) and $0$ (column $k$).  So
$\mu=(k+1,k-1,\dots,3,1)$ (i.e.\ $(4,1)$ for $k=3$) with $k-1$ parts, and $\B(\mu)$ has the parts $k-1$ and $k,k-2,\dots,2$ (the
part $1$ disappears).  That is $\B(\mu)=(k,k-1,\dots,2)=\lambda_E$ with $E=D_k\setminus\{d_k(k)\}$, which is
cyclic by Theorem~\ref{thm:brandt}.
\end{proof}

\begin{theorem}[{\cite[Thm.~4.4(2)]{GH}}]\label{thm:c3lower}
For $k\ge3$, $\lambda^{\langle k\rangle}=(k-1,k-2,k-2,k-3,\dots,2,1,1)\vdash T_k-1$ and
$\dB(\lambda^{\langle k\rangle})=k^2-2k-1$.  Consequently $\DB(T_k-1)=k^2-2k-1$ for all $k\ge4$.
\end{theorem}
\begin{proof}
The sequence is non-increasing with sum $(k-1)+(k-2)+T_{k-2}+1=T_k-1$.  Its diagram is $X'(1,k+1)$:
columns 1 and 2 have $k-1$ and $k-2$ cells (those of $\Delta_{k-1}$, so $d_k(1),d_k(2)$ are missing).  Column
$j\in[3,k]$ has $k+1-j$ cells ($\Delta_{k-1}$ plus $d_k(j)$), and column $k+1$ is $\{d_{k+1}(k+1)\}$.  Let
$h_m\equiv1+m\pmod k$ and $x_m\equiv m\pmod{k+1}$ with $h_m\in[1,k]$, $x_m\in[1,k+1]$.  Then
$(h_m,x_m)=(k-1,1)$ iff $m\equiv-2\pmod k$ and $m\equiv1\pmod{k+1}$.  $M=k^2-2k-2$ satisfies both
($k^2-2k-2\equiv1+2-2\pmod{k+1}$), and for $k\ge3$, $0\le M<k(k+1)$.  So $M$ is the least non-negative
solution.  By Lemma~\ref{lem:c3conf} and induction, $Y(\B^m(\lambda^{\langle k\rangle}))=X'(h_m,x_m)$ for $m\le M$,
and $\B^{M+1}(\lambda^{\langle k\rangle})$ is cyclic.  For $m\le M$ the diagram contains a cell of $D_{k+1}$, so
$\B^m(\lambda^{\langle k\rangle})$ is not cyclic (Corollary~\ref{cor:cyc}(b): cyclic diagrams lie in $\Delta_k$).
Hence $\dB=M+1=k^2-2k-1$.  With Theorem~\ref{thm:c3upper}, $\DB(T_k-1)=k^2-2k-1$ for $k\ge4$.  (For $k=3$,
$\DB(5)=3>2$; for $k=2$, $\DB(2)=0$.)
\end{proof}

\section{All extremal partitions of $T_k-1$}\label{sec:c3ext}

\begin{theorem}\label{thm:c3ext}
Let $k\ge5$ and $\lambda\vdash T_k-1$.  The following are equivalent:
\begin{enumerate}
\item[(1)] $\dB(\lambda)=k^2-2k-1$;
\item[(2)] $c_{k^2-2k-2}\ge k+1$, i.e.\ $\B^{k^2-2k-3}(\lambda)$ has at least $k+1$ parts;
\item[(3)] $(jk,\ jk+j+1,\ k)$ is a ramp for every $j=1,\dots,k-3$.
\end{enumerate}
\end{theorem}
\begin{proof}
(3)$\Rightarrow$(2): $j=k-3$ gives $c_{k(k-3)+k-2}=c_{k^2-2k-2}=k+1$.
(2)$\Rightarrow$(1): by Lemma~\ref{lem:overfull} (here $r=k-1$), $\dB(\lambda)\ge k^2-2k-1$; the reverse
inequality is Theorem~\ref{thm:c3upper}.
(1)$\Rightarrow$(3): go through the proof of Theorem~\ref{thm:c3upper} with $d=k^2-2k-1$.  $t\le k$ would give
$d\le k-1$.  (C) and (A) with $L=k-1$ are impossible, and whenever (B) holds it gives $d<k^2-2k-1$
(strictly, as $k\ge5$).  Hence (A) holds with $L\le k-2$, and $d\le t-1\le p+k-1\le(L-1)k+k-1\le k^2-2k-1$ is
an equality: $L=k-2$ and $p=(L-1)k$.  The equality part of Proposition~\ref{prop:rampbound} gives the ramps
$(jk,jk+j+1,k)$ for $j=1,\dots,k-3$.
\end{proof}

\begin{remark}
As for $T_k$, this is a dynamical characterisation, not a closed formula.  For $k=5$ the extremal
partitions of $14$ are $(4,3,3,2,1,1)=\lambda^{\langle5\rangle}$, $(4,3,2,2,2,1)$, $(4,3,2,2,1,1,1)$,
$(3,3,3,2,2,1)$, $(3,3,3,2,1,1,1)$ and $(3,3,2,2,2,1,1)$.  For $k=4$ the unique extremal partition of $9$ is
$(3,2,2,1,1)$; the proof of Theorem~\ref{thm:c3ext} does not apply to $k=4$, and this fact is from exhaustive
computation over the $30$ partitions of $9$ (Appendix~\ref{app:code}).  Numbers of extremal partitions of $T_k-1$:
\begin{center}
\begin{tabular}{c|ccccccc}
$k$ & 4 & 5 & 6 & 7 & 8 & 9 & 10\\\hline
$\#\{\lambda\vdash T_k-1:\dB(\lambda)=k^2-2k-1\}$ & 1 & 6 & 34 & 175 & 831 & 3911 & 18163
\end{tabular}
\end{center}
\end{remark}

\section{Towards a direct description of the extremal partitions}\label{sec:c3direct}

\begin{proposition}[a proved direct necessary condition]\label{prop:c3direct}
Let $k\ge5$ and let $\lambda\vdash T_k-1$ be extremal.  Then
\begin{enumerate}
\item[(i)] $\lambda_1\le k+1$;
\item[(ii)] $\lambda'_i:=\#\{m:\lambda_m\ge i\}\ge k+3-2i$ for every $i\ge1$.  Equivalently, $Y(\lambda)$ contains the diagram of
$\mathrm{lo}_k=(\lceil (k+2-j)/2\rceil)_{j=1}^{k+1}$, whose conjugate is $(k+1,k-1,k-3,\dots)$;
\item[(iii)] $\ell(\lambda)\ge k+1$ (the case $i=1$ of (ii)).
\end{enumerate}
\end{proposition}
\begin{proof}
By Theorem~\ref{thm:c3ext}, $(k,k+2,k)$ is a ramp: $c_k=k-1$, $c_{k+1}=k$, $c_{k+2}=k+1$.  We claim
$\R_{k+2}=\{N_1,\dots,N_{k+1}\}$.  As $c_{k+1}=c_k+1$ and $c_{k+2}=c_{k+1}+1$, Lemma~\ref{lem:basic}(c) gives
$\R_k\subseteq \R_{k+1}\subseteq \R_{k+2}$, and $N_k,N_{k+1}\in \R_{k+2}$.  Suppose some $N_i$ with $1\le i\le k-1$ is not in
$\R_{k+2}$, and take $i$ largest.  Then $N_{i+1}\in \R_{k+2}$, hence $N_{i+1}\in \R_{k+1}$ (Lemma~\ref{lem:basic}(d)), and
$N_i\notin \R_k$ (as $\R_k\subseteq \R_{k+2}$).  Lemma~\ref{lem:gap} (with $j=i<k$) gives $N_{i+1}\notin \R_{k+2}$, a contradiction.
So $N_1,\dots,N_{k+1}\in \R_{k+2}$, and since $|\R_{k+2}|=c_{k+2}=k+1$ the claim follows.

(i) No $O_m$ lies in $\R_{k+2}$, i.e.\ $\lambda_m\le k+1$ for all $m$.
(ii) For $1\le i\le k+1$, $N_i\in \R_{k+2}$ gives $c_i\ge k+2-i$.  Hence for $j<i\le k+2$, $c_j\ge k+2-j\ge i-j$, i.e.\
$N_j\in \R_i$ by \eqref{eq:alive}.  So $|\R_i|=c_i$ and \eqref{eq:alive} give $\lambda'_i=c_i-(i-1)\ge k+3-2i$.  For
$i\ge k+2$ the bound is negative.  The conjugate of $\mathrm{lo}_k$ is
$\#\{j\le k+1:\lceil(k+2-j)/2\rceil\ge i\}=\#\{j\le k+3-2i\}=k+3-2i$, which gives the equivalent form.
\end{proof}

This is the analogue of condition (1) of \cite[Thm.~3.8]{GH}.  The following observation shows why we do not
claim a direct characterisation.

\begin{remark}[computational observation, not used]\label{rem:sandwich}
Let $\mathrm{hi}_k$ be the partition with conjugate $(2k-3,2k-5,\dots,3,1)$.  For $5\le k\le10$ every extremal
partition of $T_k-1$ lies between $\mathrm{lo}_k$ and $\mathrm{hi}_k$ (\texttt{bs.py}).  But the number of partitions of
$T_k-1$ in this interval is $6,35,175,834,3935,18330$, against $6,34,175,831,3911,18163$ extremal ones.  So
``$\mathrm{lo}_k\subseteq\lambda\subseteq\mathrm{hi}_k$'' characterises the extremal partitions for $k=5,7$ but not for
$k=6,8,9,10$.  Griggs--Ho made the same kind of observation in the triangular case: their three direct
conditions \cite[Thm.~3.8]{GH} are necessary but not sufficient ($k=8$).
\end{remark}

\paragraph{What Theorem~\ref{thm:c3ext} does and does not give.}  It is an exact ``if and only if'' description of the
set of extremal partitions of $T_k-1$, and it goes beyond the bare definition in two ways.  First, (2) is a
test on a single number: the part count at one fixed time $k^2-2k-2$ must be at least $k+1$.  The definition, by
contrast, asks whether $\B^{k^2-2k-3}(\lambda)$ is cyclic, which needs Brandt's list.  Second, (3) says that every
extremal orbit has the same rigid skeleton: ramps of value $k$ at the fixed positions $jk,\dots,jk+j+1$,
$j=1,\dots,k-3$.  However, it is a \emph{dynamical} description, not a closed-form (static) description of the
diagrams.  Proposition~\ref{prop:c3direct} extracts direct conditions from the first ramp only.
Remark~\ref{rem:sandwich} shows that the natural static bounds do not suffice.  We do not have a proof that no
simple static description exists; the heuristic obstacle is that the ramps with $j\ge2$ involve the orbit up to
time about $k^2$.  \emph{We therefore regard the question ``which partitions attain the maximum'' as answered
exactly but not in closed form, and we leave a closed-form description open.}

\section{Computations}
Computations used in proofs: the case $k=4$ of Theorem~\ref{thm:c3upper} (exhaustive over the $67$ partitions
of $7,8,9$) and the uniqueness of the extremal partition of $9$ (exhaustive over its $30$ partitions).  Both are
done by the script in Appendix~\ref{app:code}, which runs in well under a second.  As sanity checks, \texttt{bs.py} (exhaustive for
$n\le60$, about $1.5$ minutes) confirms: $\DB(n)\le k^2-2k-1$ for all non-triangular $n\le60$ with $k\ge4$;
$\DB(T_k-1)=k^2-2k-1$ for $4\le k\le10$; $\dB(\lambda^{\langle k\rangle})=k^2-2k-1$ for $3\le k<80$; and
Theorem~\ref{thm:c3ext} for $5\le k\le10$.  It also checks the ramp bound and Lemma~\ref{lem:long} on
every ramp among the first $4n+10$ part counts of every partition of $n\le30$.

\appendix
\section{Code for the exhaustive small cases}\label{app:code}
File \texttt{p3\_small\_cases.py} (Python 3, no dependencies; runs in well under one second). It computes $\dB$ directly from the definition: the orbit of $\lambda$ is followed until a partition repeats, and the index of the first occurrence of that partition is $\dB(\lambda)$.
{\footnotesize
\begin{verbatim}
#!/usr/bin/env python3
"""Exhaustive small cases used in p3_c2/p3_c3/p3_c4/p3_c5.
Straight from the definitions: B(lambda), d_B(lambda) = number of steps before the
first partition that recurs in the orbit, D_B(n) = max over all partitions of n.
Runtime: well under one second."""


def partitions(n, maxp=None):
    maxp = n if maxp is None else maxp
    if n == 0:
        yield ()
        return
    for p in range(min(n, maxp), 0, -1):
        for rest in partitions(n - p, p):
            yield (p,) + rest


def B(lam):
    return tuple(sorted([len(lam)] + [x - 1 for x in lam if x > 1], reverse=True))


def d_B(lam):
    first = {}
    i = 0
    while lam not in first:
        first[lam] = i
        lam = B(lam)
        i += 1
    return first[lam]  # the tail length: B^i(lam) is cyclic iff i >= this


def fmt(lam):
    # compact notation: (3,3,2,1) -> 3^2 2 1
    out = []
    for v in sorted(set(lam), reverse=True):
        m = lam.count(v)
        out.append(f"{v}^{m}" if m > 1 else f"{v}")
    return "(" + " ".join(out) + ")"


for n in [1, 3, 4, 5, 6, 7, 8, 9, 12, 17]:
    parts = list(partitions(n))
    d = {lam: d_B(lam) for lam in parts}
    D = max(d.values())
    ext = [lam for lam in parts if d[lam] == D]
    print(f"n={n}: {len(parts)} partitions, D_B={D}, extremal: {', '.join(fmt(l) for l in ext)}")
\end{verbatim}
}
Output:
{\footnotesize
\begin{verbatim}
n=1: 1 partitions, D_B=0, extremal: (1)
n=3: 3 partitions, D_B=2, extremal: (1^3)
n=4: 5 partitions, D_B=2, extremal: (1^4)
n=5: 7 partitions, D_B=3, extremal: (1^5)
n=6: 11 partitions, D_B=6, extremal: (2^2 1^2)
n=7: 15 partitions, D_B=4, extremal: (1^7)
n=8: 22 partitions, D_B=5, extremal: (1^8)
n=9: 30 partitions, D_B=7, extremal: (3 2^2 1^2)
n=12: 77 partitions, D_B=8, extremal: (3^2 2^2 1^2), (1^12)
n=17: 297 partitions, D_B=12, extremal: (1^17)
\end{verbatim}
}

\begin{thebibliography}{9}
\bibitem{AD} E.~Akin and M.~Davis, Bulgarian solitaire, \emph{Amer.\ Math.\ Monthly} 92 (1985), no.~4, 237--250.
\bibitem{Br} J.~Brandt, Cycles of partitions, \emph{Proc.\ Amer.\ Math.\ Soc.} 85 (1982), 483--486.
\bibitem{Et} G.~Etienne, Tableaux de Young et solitaire bulgare, \emph{J.\ Combin.\ Theory Ser.~A} 58 (1991), 181--197.
\bibitem{GH} J.~R.~Griggs and C.-C.~Ho, The cycling of partitions and compositions under repeated shifts,
\emph{Adv.\ in Appl.\ Math.} 21 (1998), no.~2, 205--227. Preprint (dated March 2, 1998):
\url{https://people.math.sc.edu/griggs/cycling.pdf}. Theorem and lemma numbers refer to this preprint.
\bibitem{Ig} K.~Igusa, Solution of the Bulgarian solitaire conjecture, \emph{Math.\ Mag.} 58 (1985), 259--271.
\bibitem{OEIS} OEIS Foundation, The On-Line Encyclopedia of Integer Sequences, sequence A037306.
\end{thebibliography}

\end{document}
